\documentclass[12pt]{article}
\usepackage[utf8]{inputenc}
\usepackage[T1]{fontenc}
\usepackage{hyperref}
\title{Kumzari Language Resources for K-12 Curriculum Development}
\author{Compiled by Claude Sonnet 5.5 (claude-sonnet-5-5) for Andr\'as Kornai}
\date{}
\begin{document}
\maketitle

\section{Introduction}
Kumzari (ISO 639-3 zum; Glottocode kumz1235) is an Indo-European language of mixed Perso-Arabian character, spoken on the Musandam Peninsula of Oman and on Larak Island in Iran. It is not known to be taught in schools. Little material exists, and none of what was found is a school curriculum resource. This report lists only what appeared in search results.

\section{Primary Educational Resources}
No primers, textbooks or school materials for Kumzari were found in the searches.

\section{Linguistic and Reference Materials}
\begin{itemize}
\item A Grammar of Kumzari: a mixed Perso-Arabian language of Oman, Christina van der Wal Anonby, Leiden University doctoral dissertation (2015). \url{https://scholarlypublications.universiteitleiden.nl/access/item%3A2888893/view}
\item Study of a thematic Kumzari dictionary project (article in Arabic in the Egyptian journal JFLA). It describes the project team and its methods. A published dictionary was not confirmed. \url{https://jfla.journals.ekb.eg/article_414388.html}
\item ``The Kumzari dialect of the Shihuh tribe, Arabia, and a vocabulary'', Journal of the Royal Asiatic Society (1930). Full text is behind a paywall. \url{https://www.cambridge.org/core/journals/journal-of-the-royal-asiatic-society/article/kumzari-dialect-of-the-shihuh-tribe-arabia-and-a-vocabulary/12663D632A971AD72938E132DA7F41A7}
\item Wikipedia overview of the Kumzari language. \url{https://en.wikipedia.org/wiki/Kumzari_language}
\end{itemize}

\section{Religious and Community Resources}
None were found in the searches.

\section{Digital and Audio Resources}
\begin{itemize}
\item Bamberg WOWA metadata for the Kumzari (Musandam) dataset. It says the documentation behind Anonby's grammar is available online and includes traditional narrative texts recorded around 2010 in Kumzar. \url{https://static-export-multicast.projects.it-fopro.uni-bamberg.de/resources/wowa/data/iranian/kumzari_musandam/wowa_iran_kumzari_musandam__metadata.pdf}
\end{itemize}

\section{Curriculum Development Context}
The language is not known to be taught in schools, and the searches found no school programme. A curriculum would have to be built from the grammar, the recorded texts and the dictionary project, working with the Kumzari community. Wikipedia and the Bamberg metadata give different speaker numbers, so the figure of about 2,300 used here is uncertain.

\end{document}